% Source: https://artofproblemsolving.com/community/c3h2476592_clock_formula_challenge
% Stable identity: aops-topic-2476592-q01
\textbf{(a)} Measure positions in degrees clockwise from the top.
\begin{itemize}
\item The minute hand turns $360^\circ$ in $60$ minutes, so it is at $6m$ degrees.
\item The hour hand turns $\dfrac{360}{H}$ degrees per hour, i.e. $\dfrac{360}{60H}=\dfrac6H$ degrees per minute, so it is at $\dfrac{360h}{H}+\dfrac{6m}{H}$ degrees.
\end{itemize}
Both positions lie in $[0,360)$, because $\dfrac{360h}H+\dfrac{6m}H<\dfrac{360(H-1)}H+\dfrac{360}H=360$. The hands split the full turn into two angles: the difference of the positions,
\[\left|\dfrac{360h}{H}+\dfrac{6m}{H}-6m\right|=\left|\dfrac{360h}{H}-\dfrac{6(H-1)m}{H}\right|=\theta,\]
and $360^\circ-\theta$. The angle between the hands is the smaller one, which is $\theta$ when $\theta\le180^\circ$ and $360^\circ-\theta$ otherwise. $\blacksquare$

\textbf{(b)}
\begin{itemize}
\item $H=10$: $\theta=|36h-5.4m|$;
\item $H=6$: $\theta=|60h-5m|$;
\item $H=5$: $\theta=|72h-4.8m|$;
\item $H=24$: $\theta=|15h-5.75m|$.
\end{itemize}
(For $H=12$ this gives the familiar $|30h-5.5m|$.)

\textbf{(c)} At $7{:}30$ on a $10$-hour clock, $\theta=|252-162|=90$, so the angle is $\boxed{90^\circ}$.

At $17{:}45$ on a $24$-hour clock, $\theta=|255-258.75|=3.75$, so the angle is $\boxed{3.75^\circ}$. (Indeed the hour hand is at $17.75\cdot15=266.25^\circ$ and the minute hand at $270^\circ$.)
